\paragraph{Information Linear Conservation and Fluctuations: Mutual Information, Central Limit Theorem, and Large Deviations under Proportional Querying}
Under the microcanonical prior \(F\sim\mathrm{Unif}(\mathfrak{Fold}_m(d))\),
the type trajectory generated by any deterministic adaptive distinct-query strategy is equivalent to sampling without replacement from the multiset \(\{x^{\times d(x)}\}_{x\in X_m}\) (Theorem \ref{thm:pom-microcanonical-adaptive-no-gain-without-replacement}).
Hence, "information gain" can be characterized by purely combinatorial quantities without invoking geometric details: the logarithmic contraction of the posterior moduli space \(\mathcal{M}_k\) yields the information density;
its thermodynamic limit is linear on the exponential scale and exhibits Gaussian fluctuations on the \(\sqrt N\) scale; meanwhile a full large deviation description is provided by the multivariate hypergeometric KL-symmetric rate function.

\begin{definition}[Information Density and Normalized Information Rate]\label{def:pom-microcanonical-information-density}
In the setting of Definitions \ref{def:pom-microcanonical-adaptive-trajectory-counts} and \ref{def:pom-microcanonical-posterior-moduli-size-beta},
define
\[
\mathsf{I}_k:=\log\card{\mathfrak{Fold}_m(d)}-\log \mathcal{M}_k(T_{1:k})
\]
as the \emph{information density} (that is, the self-information of the observed trajectory, see Theorem \ref{thm:pom-microcanonical-information-identity}).
Define the normalized quantities
\[
\iota_N:=\frac{1}{N}\mathsf{I}_k,
\qquad
\ell_N:=\frac{1}{N}\log \mathcal{M}_k(T_{1:k}).
\]
\end{definition}

\begin{theorem}[Exact Moduli Space Representation of Information Density]\label{thm:pom-microcanonical-information-identity}
Under the above microcanonical prior and deterministic adaptive distinct-query setting, for any \(k\le N\) there holds the strict identity
\[
\boxed{\;
\mathsf{I}_k
=-\log\mathbb{P}(T_{1:k})
=\log\card{\mathfrak{Fold}_m(d)}-\log \mathcal{M}_k(T_{1:k})\;}
\]
and the mutual information satisfies
\[
\boxed{\;
I(F;T_{1:k})=H(T_{1:k})
=\EE[\mathsf{I}_k]
=\log\card{\mathfrak{Fold}_m(d)}-\EE\bigl[\log \mathcal{M}_k(T_{1:k})\bigr]\;}
\]
where \(H(T_{1:k})\) is the Shannon entropy (natural logarithm) of the trajectory distribution.
\end{theorem}
\begin{proof}
Because the query strategy is deterministic and each query point is uniquely determined by the historical responses, the trajectory \(T_{1:k}\) is a deterministic function of \(F\), hence \(I(F;T_{1:k})=H(T_{1:k})\).
\par
On the other hand, by the trajectory probability formula of Theorem \ref{thm:pom-microcanonical-fold-posterior-count-and-prob},
for any realizable trajectory \(t\in X_m^k\), its probability is exactly
\[
\mathbb{P}(T_{1:k}=t)=\frac{\card{\{F'\in\mathfrak{Fold}_m(d):F'(\omega_i)=t_i,\ 1\le i\le k\}}}{\card{\mathfrak{Fold}_m(d)}}
=\frac{\mathcal{M}_k(t)}{\card{\mathfrak{Fold}_m(d)}}.
\]
Taking \(t=T_{1:k}\) yields
\(
-\log\mathbb{P}(T_{1:k})
=\log\card{\mathfrak{Fold}_m(d)}-\log\mathcal{M}_k(T_{1:k})
\), establishing the first equation.
Taking expectations on both sides and noting \(\EE[-\log\mathbb{P}(T_{1:k})]=H(T_{1:k})\) yields the second set of identities. This completes the proof.
\end{proof}

\begin{theorem}[Information Linear Conservation Law: Exponential Rate under Proportional Querying]\label{thm:pom-microcanonical-information-linear-law}
Take constant \(\beta\in[0,1]\) and set \(k=\lfloor\beta N\rfloor\).
Then as \(N\to\infty\) with fixed \(w\), we have convergence in probability
\[
\boxed{\;
\frac{1}{N}\mathsf{I}_k\ \xrightarrow{\ \mathbb{P}\ }\ \beta\,H(w)\;},
\qquad
\boxed{\;
\frac{1}{k}\mathsf{I}_k\ \xrightarrow{\ \mathbb{P}\ }\ H(w)\;}
\quad(\beta\in(0,1]),
\]
where \(w(x)=d(x)/N\) and \(H(w)=-\sum_x w(x)\log w(x)\).
Moreover, the mutual information satisfies the same exponential rate
\[
\frac{1}{N}I(F;T_{1:k})=\frac{1}{N}\EE[\mathsf{I}_k]\ \to\ \beta\,H(w).
\]
\end{theorem}
\begin{proof}
By Theorem \ref{thm:pom-microcanonical-information-identity},
\(
\mathsf{I}_k=\log\card{\mathfrak{Fold}_m(d)}-\log\mathcal{M}_k(T_{1:k})
\).
By Theorem \ref{thm:pom-microcanonical-fold-entropy},
\[
\frac{1}{N}\log\card{\mathfrak{Fold}_m(d)}=H(w)+o(1),
\]
while by Theorem \ref{thm:pom-microcanonical-posterior-entropy-linear-law} we have
\[
\frac{1}{N}\log\mathcal{M}_k(T_{1:k})\ \xrightarrow{\ \mathbb{P}\ }\ (1-\beta)\,H(w).
\]
Subtracting these two equations yields \(\frac{1}{N}\mathsf{I}_k\to \beta H(w)\) in probability.
Dividing by \(k/N\to\beta\) yields the second formula.
The last formula follows from \(\EE[\mathsf{I}_k]=I(F;T_{1:k})\) and uniform boundedness of \(\mathsf{I}_k/N\): since \(\log\mathcal{M}_k(T_{1:k})\ge 0\), we have
\[
0\le \frac{1}{N}\mathsf{I}_k\le \frac{1}{N}\log\card{\mathfrak{Fold}_m(d)}=H(w)+o(1),
\]
hence \(\mathsf{I}_k/N\) is uniformly bounded in \(L^1\); combined with convergence in probability this gives
\(
\EE[\mathsf{I}_k]/N\to \beta H(w)
\).
This completes the proof.
\end{proof}

\begin{corollary}[Central Limit Theorem for Mutual Information Fluctuations]\label{cor:pom-microcanonical-information-clt}
In the setting of Theorem \ref{thm:pom-microcanonical-posterior-moduli-clt} (in particular with fixed \(n\), fixed \(w\), and \(\beta\in(0,1)\)),
write
\[
V(w):=\Var_{X\sim w}\bigl(\log w(X)\bigr).
\]
Then as \(N\to\infty\),
\[
\boxed{\;
\frac{\mathsf{I}_k-\beta N H(w)}{\sqrt{\beta(1-\beta)N}}
\ \xrightarrow{d}\ \mathcal N\bigl(0,V(w)\bigr)\;}
\qquad\left(k=\lfloor\beta N\rfloor\right).
\]
\end{corollary}
\begin{proof}
By Definition \ref{def:pom-microcanonical-information-density} and Theorem \ref{thm:pom-microcanonical-fold-entropy},
\[
\mathsf{I}_k
=\log\card{\mathfrak{Fold}_m(d)}-\log\mathcal{M}_k(T_{1:k})
=NH(w)-\log\mathcal{M}_k(T_{1:k})+O(n\log N).
\]
Substituting the central limit theorem of Theorem \ref{thm:pom-microcanonical-posterior-moduli-clt} and noting that \(O(n\log N)=o(\sqrt N)\) when \(n\) is fixed, we obtain the conclusion. This completes the proof.
\end{proof}

\begin{theorem}[Large Deviations for Multivariate Hypergeometric: KL-Symmetric Decomposition of Sample and Residual Compositions]\label{thm:pom-microcanonical-hypergeometric-ldp-kl-sym}
Assume fixed \(n\), fixed \(w\), and \(\beta\in(0,1)\).
Let \(k=\lfloor\beta N\rfloor\), and define sample composition and residual composition
\[
Q_N(x):=\frac{C_k(x)}{k}\in\Delta(X_m),
\qquad
R_N(x):=\frac{d(x)-C_k(x)}{N-k}\in\Delta(X_m).
\]
Then \(Q_N\) satisfies a large deviation principle with good rate function given explicitly by
\[
J_\beta(q)=
\begin{cases}
\beta\,D(q\|w)+(1-\beta)\,D(r(q)\|w), & r(q):=\dfrac{w-\beta q}{1-\beta}\in\Delta(X_m),\\[6pt]
+\infty, & \text{otherwise},
\end{cases}
\]
where \(D(\cdot\|\cdot)\) is the KL divergence (natural logarithm).
Moreover, the symmetry property holds
\[
\boxed{\;
J_\beta(q)=J_{1-\beta}\bigl(r(q)\bigr)\;}
\]
corresponding to the exchange invariance of mass conservation \(w=\beta q+(1-\beta)r\).
\end{theorem}
\begin{proof}
By the probability mass function of the hypergeometric distribution (the counting form of Theorem \ref{thm:pom-microcanonical-adaptive-no-gain-without-replacement}),
for integer vector \(c\) satisfying \(\sum_x c(x)=k\) we have
\[
\mathbb{P}(C_k=c)=\frac{\prod_{x\in X_m}\binom{d(x)}{c(x)}}{\binom{N}{k}}.
\]
Taking \(c(x)=kq(x)\) and expanding the factorials using Stirling's formula, we obtain
\[
-\frac1N\log \mathbb{P}(Q_N\approx q)
=\beta\,D(q\|w)+(1-\beta)\,D(r(q)\|w)+o(1),
\]
where \(r(q)=(w-\beta q)/(1-\beta)\) is forced by the conservation relation.
This expansion holds when \(q\) is an interior point of the simplex and \(r(q)\in\Delta(X_m)\), yielding the good rate function; if \(r(q)\notin\Delta\) then the probability is \(0\) and the rate is \(+\infty\).
Symmetry follows directly from exchanging \(\beta q\) and \((1-\beta)r\) while keeping \(w\) invariant. This completes the proof.
\end{proof}

\begin{theorem}[Full Large Deviations for Posterior Moduli Space and Information Density: Explicit Variational Form via Contraction Principle]\label{thm:pom-microcanonical-moduli-information-ldp}
Under the assumptions of Theorem \ref{thm:pom-microcanonical-hypergeometric-ldp-kl-sym}, both \(\ell_N\) and \(\iota_N\) satisfy large deviation principles.
Their rate functions can be written as
\[
K_\ell(s)=\inf\Bigl\{J_\beta(q):\ s=(1-\beta)\,H(r(q))\Bigr\},
\]
\[
K_\iota(s)=\inf\Bigl\{J_\beta(q):\ s=H(w)-(1-\beta)\,H(r(q))\Bigr\},
\]
where \(r(q)=(w-\beta q)/(1-\beta)\), and \(H(\cdot)\) is Shannon entropy (natural logarithm).
The minimum is attained at \(q=w\) (hence \(r(w)=w\)), yielding the typical values
\[
\ell_N\ \to\ (1-\beta)H(w),
\qquad
\iota_N\ \to\ \beta H(w),
\]
consistent with Theorem \ref{thm:pom-microcanonical-information-linear-law}.
\end{theorem}
\begin{proof}
On the event \(\{Q_N\approx q\}\), the residual composition satisfies \(R_N\approx r(q)\).
By Stirling expansion (see the analogous derivation in Theorem \ref{thm:pom-microcanonical-posterior-entropy-linear-law}) we obtain
\[
\frac{1}{N}\log\mathcal{M}_k(T_{1:k})
=(1-\beta)H(R_N)+o(1)
=(1-\beta)H(r(q))+o(1).
\]
Hence within the effective domain, the map \(q\mapsto (1-\beta)H(r(q))\) is continuous.
Applying the contraction principle to the large deviation principle for \(Q_N\) yields \(K_\ell\).
By Theorem \ref{thm:pom-microcanonical-fold-entropy}, we have \(\frac{1}{N}\log\card{\mathfrak{Fold}_m(d)}=H(w)+o(1)\).
Therefore
\[
\iota_N=\frac{1}{N}\log\card{\mathfrak{Fold}_m(d)}-\ell_N
=H(w)-\ell_N+o(1),
\]
and deterministic \(o(1)\) translation does not affect the large deviation rate function on the \(N\) scale;
hence \(K_\iota\) follows immediately from \(K_\ell\) (can also be viewed as the equivalent form \(K_\iota(s)=K_\ell(H(w)-s)\)).
The minimum is obtained from strict convexity of \(J_\beta\) (in the interior) and \(J_\beta(w)=0\). This completes the proof.
\end{proof}

\begin{theorem}[Doob Decomposition of Entropy and Information: Conditional Increments and Non-Asymptotic Concentration]\label{thm:pom-microcanonical-information-doob}
For any \(i\in\{0,1,\dots,k\}\) define residual composition
\[
p_i(x):=\frac{d(x)-C_i(x)}{N-i}\in\Delta(X_m),
\qquad C_0\equiv 0,
\]
and let filtration \(\mathcal F_i:=\sigma(T_1,\dots,T_i)\).
Then the following holds:
\begin{enumerate}
  \item \textbf{(Conditional Increment)} For each \(1\le i\le k\),
  \[
  \mathbb{P}(T_i=x\mid \mathcal F_{i-1})=p_{i-1}(x),
  \qquad
  \mathsf{I}_k=\sum_{i=1}^k\bigl(-\log p_{i-1}(T_i)\bigr)\ \ \text{a.s.}
  \]
  and
  \[
  \EE\!\left[-\log\mathbb{P}(T_i\mid\mathcal F_{i-1})\ \big|\ \mathcal F_{i-1}\right]=H(p_{i-1}).
  \]
  \item \textbf{(Martingale Structure)} \((p_i)_{i\ge 0}\) is a \(\Delta(X_m)\)-valued martingale: \(\EE[p_i\mid\mathcal F_{i-1}]=p_{i-1}\).
  \item \textbf{(Azuma-Type Concentration)} Let
  \[
  M_k:=\mathsf{I}_k-\sum_{i=1}^k H(p_{i-1}),
  \]
  then \((M_i)\) is a martingale with increments satisfying \(|M_i-M_{i-1}|\le \log N\).
  Hence for any \(t>0\) we have the non-asymptotic bound
  \[
  \boxed{\;
  \mathbb{P}\Bigl(\bigl|\mathsf{I}_k-\sum_{i=1}^k H(p_{i-1})\bigr|\ge t\Bigr)
  \le 2\exp\!\Bigl(-\frac{2t^2}{k(\log N)^2}\Bigr)\; } .
  \]
\end{enumerate}
\end{theorem}
\begin{proof}
Point (1) follows directly from the stepwise conditional distribution formula of Theorem \ref{thm:pom-microcanonical-adaptive-no-gain-without-replacement}.
By chain rule,
\(
\mathbb{P}(T_{1:k})=\prod_{i=1}^k\mathbb{P}(T_i\mid\mathcal F_{i-1})
\), taking negative logarithm yields \(\mathsf{I}_k=\sum_i -\log p_{i-1}(T_i)\).
The conditional expectation identity follows from
\(
\EE[-\log p_{i-1}(T_i)\mid\mathcal F_{i-1}]
=\sum_x p_{i-1}(x)(-\log p_{i-1}(x))
=H(p_{i-1})
\).
\par
Point (2): For any \(x\), from \(C_i(x)=C_{i-1}(x)+\mathbf{1}_{\{T_i=x\}}\) and
\(\EE[\mathbf{1}_{\{T_i=x\}}\mid\mathcal F_{i-1}]=p_{i-1}(x)\) we obtain
\[
\EE[p_i(x)\mid\mathcal F_{i-1}]
=\EE\!\left[\frac{d(x)-C_{i-1}(x)-\mathbf{1}_{\{T_i=x\}}}{N-i}\ \Big|\ \mathcal F_{i-1}\right]
=\frac{(N-i+1)p_{i-1}(x)-p_{i-1}(x)}{N-i}
=p_{i-1}(x).
\]
\par
Point (3): By point (1), \(M_i-M_{i-1}=-\log p_{i-1}(T_i)-H(p_{i-1})\) has conditional mean \(0\), hence \(M_i\) is a martingale.
Moreover, since \(p_{i-1}(T_i)\ge 1/(N-i+1)\) (the observed type appears at least once in the residual multiset),
we have \(-\log p_{i-1}(T_i)\le \log(N-i+1)\le \log N\), while \(0\le H(p_{i-1})\le \log n\le \log N\);
thus \(|M_i-M_{i-1}|\le \log N\).
Applying Azuma-Hoeffding to \((M_i)\) yields the stated bound. This completes the proof.
\end{proof}

\begin{theorem}[Query + Side Information Bits: Optimal Recovery Rate and Gaussian Critical Window]\label{thm:pom-microcanonical-query-plus-bits-phase}
After observing \(k\) query trajectories \(T_{1:k}\), allow additional external information of \(B\) bits:
the encoder can take any map \(\mathrm{enc}:\mathfrak{Fold}_m(d)\to\{0,1\}^B\), and the decoder outputs \(\widehat F\) given input \((T_{1:k},\mathrm{enc}(F))\).
\begin{enumerate}
  \item \textbf{(Optimal Success Rate Closed Form)} Under uniform prior, the optimal exact recovery success rate satisfies the strict identity
  \[
  \boxed{\;
  \mathrm{Succ}^{\mathrm{opt}}_{k,B}
  =\EE\!\left[\min\!\left(1,\frac{2^B}{\mathcal{M}_k(T_{1:k})}\right)\right]\;}
  \]
  where \(\mathcal{M}_k(T_{1:k})\) is the posterior compatible set size (Definition \ref{def:pom-microcanonical-posterior-moduli-size-beta}).
  \item \textbf{(Threshold and Gaussian Window)}
  Assume fixed \(n\), fixed \(w\), and \(\beta\in(0,1)\), let \(k=\lfloor\beta N\rfloor\) and take
  \[
  B_N=\frac{(1-\beta)NH(w)}{\log 2}
  +\frac{t}{\log 2}\sqrt{\beta(1-\beta)N V(w)},
  \qquad t\in\RR.
  \]
  Then as \(N\to\infty\),
  \[
  \boxed{\;
  \mathrm{Succ}^{\mathrm{opt}}_{k,B_N}\ \longrightarrow\ \Phi(t)\;}
  \]
  where \(\Phi\) is the standard normal distribution function.
  In particular, if \(B<(1-\beta)NH(w)/\log 2-\omega(\sqrt N)\), then \(\mathrm{Succ}^{\mathrm{opt}}_{k,B}\to 0\);
  if \(B>(1-\beta)NH(w)/\log 2+\omega(\sqrt N)\), then \(\mathrm{Succ}^{\mathrm{opt}}_{k,B}\to 1\).
\end{enumerate}
\end{theorem}
\begin{proof}
Point (1): Conditioned on trajectory \(T_{1:k}\), the posterior is uniform on the compatible set (Theorem \ref{thm:pom-microcanonical-count-sufficient-statistic-posterior-uniform}), whose size is \(\mathcal{M}_k(T_{1:k})\).
Any \(B\)-bit encoding can distinguish at most \(2^B\) candidate maps; conditioned on the trajectory, the optimal procedure is to select at most \(2^B\) maps within the posterior block and assign distinct codewords, merging the remaining maps to any codeword.
Hence the conditional success rate's optimal value is \(\min(1,2^B/\mathcal{M}_k)\); taking expectation over trajectories yields the result.
\par
Point (2): By Theorem \ref{thm:pom-microcanonical-posterior-moduli-clt},
\[
\frac{\log \mathcal{M}_k-(1-\beta)NH(w)}{\sqrt{\beta(1-\beta)N}}
\ \xrightarrow{d}\ \mathcal N\bigl(0,V(w)\bigr).
\]
Let \(B_N\log 2=(1-\beta)NH(w)+t\sqrt{\beta(1-\beta)NV(w)}\).
Then
\[
\min\!\left(1,\frac{2^{B_N}}{\mathcal{M}_k}\right)
=\min\!\left(1,\exp\!\bigl(B_N\log 2-\log\mathcal{M}_k\bigr)\right)
\xrightarrow{d}\ \mathbf{1}_{\{Z\le t\}},
\]
where \(Z\sim\mathcal N(0,1)\).
Since the left side always lies in \([0,1]\), by dominated convergence (or Portmanteau + boundedness) taking expectation yields
\(
\mathrm{Succ}^{\mathrm{opt}}_{k,B_N}\to \mathbb{P}(Z\le t)=\Phi(t)
\).
The two-sided sharp transitions follow from \(t\to\pm\infty\). This completes the proof.
\end{proof}
% (User input window)
% 2026-02-14 11:55:10 (Asia/Singapore)

\paragraph{Exponential Volume of Type-Class Hamming Balls and Partial Inversion Threshold: Diagonal-Constrained Maximum Entropy Coupling}
The microcanonical posterior exhibits typical type-class structure when conditioned on "fixing only type counts/marginal composition": given marginal distribution \(w\), unlabeled vectors are uniform on the length-\(t\) type class.
Therefore, the optimal "partial inversion" success rate under \(\delta\)-distortion can be reduced to a purely combinatorial covering problem:
how many intra-type Hamming balls are needed to cover the type class on the exponential scale.
This paragraph provides the exact exponential threshold of this covering rate, representing it as a maximum entropy coupling problem under diagonal constraints;
further proves uniqueness of the maximal coupling, exponential family structure, scalarization solution, and mutual information/dual derivative characterization.

\begin{definition}[Type Class and Intra-Type Hamming Ball]\label{def:pom-typeclass-and-hamming-ball}
Let \(X\) be a finite alphabet with \(|X|=n\ge 2\), take full-support distribution \(w\in\Delta(X)\).
Let \(t\to\infty\) satisfying \(t\,w(x)\in\NN\) (for all \(x\in X\)).
Define the type class
\[
\mathcal T_t(w):=\Bigl\{x^t\in X^t:\ \frac{1}{t}\bigl|\{i:\ x_i=a\}\bigr|=w(a)\ \ \forall a\in X\Bigr\}.
\]
For any reference sequence \(x^t\in\mathcal T_t(w)\) and any distortion parameter \(\delta\in[0,1]\), define the intra-type ball
\[
\mathcal B_t(w,\delta;x^t):=\Bigl\{y^t\in\mathcal T_t(w):\ t^{-1}d_{\rm H}(x^t,y^t)\le \delta\Bigr\}.
\]
\end{definition}

\begin{definition}[Exponential Volume and Covering Threshold]\label{def:pom-typeclass-volume-and-threshold}
If the limit exists and is independent of the reference sequence \(x^t\in\mathcal T_t(w)\), define the exponential volume
\[
\mathcal V_w(\delta):=\lim_{t\to\infty}\frac{1}{t}\log|\mathcal B_t(w,\delta;x^t)|.
\]
Define the corresponding covering threshold
\[
R_w(\delta):=H(w)-\mathcal V_w(\delta)\ \ge\ 0.
\]
\end{definition}

\begin{theorem}[Existence and Variational Representation of Exponential Volume]\label{thm:pom-typeclass-hamming-ball-volume-variational}
In the setting of Definition \ref{def:pom-typeclass-and-hamming-ball}, the limit \(\mathcal V_w(\delta)\) exists and satisfies the strict variational formula
\[
\boxed{\;
\mathcal V_w(\delta)
=
\sup\Bigl\{H_P(Y|X):\ P\in\Delta(X\times X),\ P_X=P_Y=w,\ \mathbb P_P(X=Y)\ge 1-\delta\Bigr\}\; }.
\]
Moreover, this exponential volume is independent of the choice of reference sequence \(x^t\in\mathcal T_t(w)\).
\end{theorem}
\begin{proof}
For any \(x^t\in\mathcal T_t(w)\) and any \(y^t\in\mathcal T_t(w)\), denote their joint empirical distribution (contingency table) as \(P_{x^t y^t}\in\Delta(X\times X)\).
The type constraint \(y^t\in\mathcal T_t(w)\) is equivalent to \(P_Y=w\), while \(x^t\in\mathcal T_t(w)\) is equivalent to \(P_X=w\).
The Hamming constraint \(t^{-1}d_{\rm H}(x^t,y^t)\le \delta\) is equivalent to \(P(X=Y)\ge 1-\delta\).
\par
For fixed \(x^t\), when \(y^t\) has joint type equal to some feasible \(P\),
the number of \(y^t\) conforming to that joint type has exponential rate \(H_P(Y|X)\) (standard type counting: \(|\{y^t:\ P_{x^t y^t}=P\}|=\exp(tH_P(Y|X)+o(t))\)).
Hence \(|\mathcal B_t(w,\delta;x^t)|\) is determined on the exponential scale by the maximum \(H_P(Y|X)\) over the feasible set, yielding the stated supremum.
\par
Since the feasible set is compact in the finite-dimensional simplex and \(H_P(Y|X)\) is continuous, the maximum exists.
Upper bound (\(\limsup\)) and lower bound (\(\liminf\)) are obtained by pairing upper bounds over unions of finitely many joint types with lower bounds by selecting near-optimal joint types, establishing existence of the limit.
Finally, the type class has transitive action under coordinate permutations, \(|\mathcal B_t(w,\delta;x^t)|\) depends only on the type of \(x^t\) rather than the specific sequence, hence is independent of the reference sequence. This completes the proof.
\end{proof}

\begin{theorem}[Unique Maximal Coupling and Exponential Family Structure]\label{thm:pom-typeclass-hamming-ball-unique-maximizer-exponential-family}
For each \(\delta\in(0,1)\), the supremum in Theorem \ref{thm:pom-typeclass-hamming-ball-volume-variational} is attained by a unique coupling \(P^\star_\delta\).
There exist unique scalar \(\lambda(\delta)\in\RR\) and unique positive vector \(u(\delta)\in\RR_{>0}^{X}\) such that
\[
\boxed{\;
P^\star_\delta(x,y)
=
\frac{u_xu_y\,\exp\!\bigl(\lambda(\delta)\,\mathbf 1_{\{x=y\}}\bigr)}
{\sum_{a,b\in X}u_au_b\,\exp\!\bigl(\lambda(\delta)\,\mathbf 1_{\{a=b\}}\bigr)}\qquad(\forall x,y\in X)\; }.
\]
Moreover, its marginals satisfy \(P^\star_{\delta,X}=P^\star_{\delta,Y}=w\).
Write \(p_0:=\sum_{x\in X}w(x)^2\).
If \(1-\delta>p_0\), the diagonal mass constraint saturates and satisfies \(\mathbb P_{P^\star_\delta}(X=Y)=1-\delta\), with \(\lambda(\delta)>0\).
If \(1-\delta\le p_0\), the optimal coupling is the independent coupling \(P^\star_\delta=w\otimes w\), and \(\lambda(\delta)=0\), \(\mathcal V_w(\delta)=H(w)\), \(R_w(\delta)=0\).
\end{theorem}
\begin{proof}
The feasible set is given by linear constraints and inequality constraints, hence is compact and convex.
The function \(H_P(Y|X)=H(P)-H(w)\) is strictly concave in \(P\) (since \(H(P)\) is strictly concave and \(H(w)\) is constant),
hence the maximizer is unique.
\par
Duality (KKT) yields exponential family form:
\(
P(x,y)\propto \exp(\alpha_x+\beta_y+\lambda\mathbf 1_{\{x=y\}})
\).
Since constraints require \(P_X=P_Y=w\) and the problem is symmetric under exchanging \(X,Y\), uniqueness forces \(\alpha=\beta\) (up to constant absorbed in normalization).
Setting \(u_x=e^{\alpha_x}>0\) yields the stated structure.
\par
Let \(P^{\mathrm{ind}}:=w\otimes w\), then \(H_{P^{\mathrm{ind}}}(Y|X)=H(w)\) and \(\mathbb P_{P^{\mathrm{ind}}}(X=Y)=p_0\).
Since for any coupling \(P\) we always have \(H_P(Y|X)\le H(Y)=H(w)\), when \(1-\delta\le p_0\), \(P^{\mathrm{ind}}\) is feasible and attains the global upper bound;
hence by uniqueness we know \(P^\star_\delta=P^{\mathrm{ind}}\), and the exponential family parameter satisfies \(\lambda(\delta)=0\).
\par
When \(1-\delta>p_0\), the independent coupling is infeasible.
If the constraint is strict inequality at the optimal point, then by complementary slackness we must have \(\lambda(\delta)=0\), yielding the independent coupling again, a contradiction.
Hence the constraint must saturate, with \(\lambda(\delta)>0\) and \(\mathbb P_{P^\star_\delta}(X=Y)=1-\delta\). This completes the proof.
\end{proof}

\begin{theorem}[Exponential Threshold for Partial Inversion: Covering Rate Determined by \texorpdfstring{$R_w(\delta)$}{Rw(delta)}]\label{thm:pom-microcanonical-partial-inversion-threshold}
Let \(Y^t\) be uniformly distributed on \(\mathcal T_t(w)\).
Allow additional \(B\) bits of side information: encoder takes any map \(\mathrm{enc}:\mathcal T_t(w)\to\{1,\dots,2^B\}\),
decoder outputs \(\widehat Y^t\in\mathcal T_t(w)\) given input \(\mathrm{enc}(Y^t)\).
For fixed \(\delta\in(0,1)\) define distortion event
\[
\mathcal E_\delta:=\{t^{-1}d_{\rm H}(Y^t,\widehat Y^t)\le \delta\},
\]
and let optimal success rate be
\[
\mathrm{Succ}^{\mathrm{opt}}_t(B,\delta):=\sup_{\mathrm{enc},\mathrm{dec}}\ \mathbb P(\mathcal E_\delta).
\]
Let \(\mathcal V_w(\delta)\) and \(R_w(\delta)\) be as in Definition \ref{def:pom-typeclass-volume-and-threshold}.
As \(t\to\infty\) with \(B=\lfloor rt/\log 2\rfloor\) (rate \(r\ge 0\) fixed), a phase transition occurs:
\[
\boxed{\;
\lim_{t\to\infty}\mathrm{Succ}^{\mathrm{opt}}_t(B,\delta)=
\begin{cases}
0, & r<R_w(\delta),\\
1, & r>R_w(\delta).
\end{cases}\;}
\]
Moreover, when \(r<R_w(\delta)\) there is an exponential upper bound
\[
\boxed{\;
\limsup_{t\to\infty}\frac{1}{t}\log \mathrm{Succ}^{\mathrm{opt}}_t(B,\delta)\ \le\ r-R_w(\delta)\;<0\; }.
\]
\end{theorem}
\begin{proof}
By type counting, \(|\mathcal T_t(w)|=\exp(tH(w)+o(t))\).
By Theorem \ref{thm:pom-typeclass-hamming-ball-volume-variational}, for any reference point \(x^t\in\mathcal T_t(w)\) we have
\(|\mathcal B_t(w,\delta;x^t)|=\exp(t\mathcal V_w(\delta)+o(t))\), with error term independent of reference point.
\par
\textbf{Upper bound:}
Any encoding produces at most \(M=2^B=\exp(rt+o(t))\) decoding output points,
whose \(\delta\)-success set is contained in the union of at most \(M\) intra-type balls, with covering number at most
\(
M\exp(t\mathcal V_w(\delta)+o(t))
\).
Dividing by \(|\mathcal T_t(w)|\) yields
\[
\mathrm{Succ}^{\mathrm{opt}}_t(B,\delta)
\le
\exp\!\bigl(t(r-\!R_w(\delta))+o(t)\bigr),
\]
establishing exponential decay when \(r<R_w(\delta)\) and yielding the stated \(\limsup\) upper bound.
\par
\textbf{Lower bound:}
Take a maximal ball covering (greedy selection of ball centers suffices), the required number of balls has exponential rate equal to
\(|\mathcal T_t(w)|/|\mathcal B_t(w,\delta;\cdot)|=\exp(tR_w(\delta)+o(t))\).
When \(r>R_w(\delta)\), the allowed number of codewords \(M=\exp(rt+o(t))\) strictly exceeds the required covering number on the exponential scale,
hence one can construct a codebook with coverage ratio tending to \(1\), yielding \(\mathrm{Succ}^{\mathrm{opt}}_t(B,\delta)\to 1\).
This completes the proof.
\end{proof}

\paragraph{Scalarization of Optimal Coupling: Single-Parameter Solution for Diagonal Multiplier}
Following the exponential family parameters of Theorem \ref{thm:pom-typeclass-hamming-ball-unique-maximizer-exponential-family}.
Write \(\kappa:=e^{\lambda}-1\in(-1,\infty)\), and denote \(A:=\sum_{x\in X}u_x\).

\begin{theorem}[Quadratic Roots and Self-Consistent Equation for Marginal Constraints]\label{thm:pom-typeclass-diagonal-coupling-scalarization}
For any given \(\kappa>-1\) and any \(A>0\), the marginal constraint \(P_X=P_Y=w\) is equivalent to satisfying the quadratic equation for each \(x\in X\)
\[
\boxed{\;
\kappa u_x^2+A u_x-w(x)=0\; } .
\]
Hence (taking positive root) we have explicit expression
\[
u_x(A,\kappa)=
\begin{cases}
\dfrac{-A+\sqrt{A^2+4\kappa w(x)}}{2\kappa}, & \kappa\ne0,\\[10pt]
\dfrac{w(x)}{A}, & \kappa=0.
\end{cases}
\]
Moreover, for each fixed \(\kappa>-1\), there exists unique \(A=A(\kappa)>0\) such that the self-consistent condition \(A=\sum_x u_x(A,\kappa)\) holds;
this unique solution yields unique vector \(u(\kappa)\).
\end{theorem}
\begin{proof}
By the exponential family form we can write (where \(A=\sum_x u_x\))
\[
P(x,y)=\frac{u_xu_y\bigl(1+\kappa\mathbf 1_{\{x=y\}}\bigr)}{Z},
\qquad
Z:=\sum_{a,b\in X}u_au_b\bigl(1+\kappa\mathbf 1_{\{a=b\}}\bigr)=A^2+\kappa\sum_{a\in X}u_a^2.
\]
Summing over fixed \(x\) yields marginal
\[
w(x)=\sum_y P(x,y)=\frac{u_x(A+\kappa u_x)}{Z}.
\]
Note that \(u\mapsto cu\ (c>0)\) does not change \(P\) (numerator and denominator multiply by \(c^2\)), so we can choose equivalent scaling such that \(Z=1\).
Under this normalization, the marginal constraint is equivalent to \(w(x)=u_x(A+\kappa u_x)\), i.e., the stated quadratic equation.
The explicit positive root follows immediately.
\par
For fixed \(\kappa>-1\), the function \(A\mapsto \sum_x u_x(A,\kappa)\) is continuous and strictly decreasing, tending to \(+\infty\) as \(A\downarrow 0\) and to \(0\) as \(A\to\infty\),
hence the equation \(A=\sum_x u_x(A,\kappa)\) has a unique solution.
Uniqueness yields uniqueness of \(u(\kappa)\).
Finally, summing \(w(x)=u_x(A+\kappa u_x)\) over \(x\) gives \(A^2+\kappa\sum_x u_x^2=\sum_x w(x)=1\), hence \(Z=1\) is consistent with the above normalization. This completes the proof.
\end{proof}

\begin{theorem}[Strict Monotonicity and Unique Matching of Diagonal Mass]\label{thm:pom-typeclass-diagonal-mass-monotone}
Let \(P^\star_{\delta(\kappa)}\) be the optimal coupling given by parameter \(\kappa\in(-1,\infty)\), and denote its diagonal mass
\[
\Delta(\kappa):=\sum_{x\in X}P^\star_{\delta(\kappa)}(x,x)=\mathbb P(X=Y).
\]
Then \(\Delta(\kappa)\) is strictly increasing on \(\kappa\in(-1,\infty)\), with endpoint values
\[
\Delta(0)=\sum_x w(x)^2,\qquad
\lim_{\kappa\downarrow-1}\Delta(\kappa)=0,\qquad
\lim_{\kappa\to\infty}\Delta(\kappa)=1.
\]
Hence when \(1-\delta>\sum_x w(x)^2\), there exists unique \(\kappa(\delta)>0\) such that \(\Delta(\kappa(\delta))=1-\delta\);
while when \(1-\delta\le \sum_x w(x)^2\), the optimal coupling is the independent coupling corresponding to \(\kappa(\delta)=0\).
\end{theorem}
\begin{proof}
By the exponential family form and marginal constraint, \(\Delta(\kappa)\) can be expressed as a rational function of \(u(\kappa)\);
when \(\kappa\) increases, the diagonal weight \((1+\kappa)u_x^2\) strictly rises relative to off-diagonal weight \(u_xu_y\ (x\neq y)\),
while the marginal is fixed to \(w\), hence diagonal mass must strictly increase.
Endpoint values follow from continuity as \(\kappa\to 0\) and limits as \(\kappa\downarrow-1,\ \kappa\to\infty\) (diagonal terms vanish/dominate respectively).
Strict monotonicity and endpoint surjectivity guarantee: for any target diagonal mass \(p\in(0,1)\) there exists unique \(\kappa\in(-1,\infty)\) such that \(\Delta(\kappa)=p\).
Combined with the discussion in Theorem \ref{thm:pom-typeclass-hamming-ball-unique-maximizer-exponential-family} of "active constraint \(\Leftrightarrow 1-\delta>\sum_x w(x)^2\)", we obtain the stated matching relation. This completes the proof.
\end{proof}

\begin{theorem}[Keep-Resample Structure: Optimal Coupling as One-Parameter Gibbs Channel]\label{thm:pom-typeclass-keep-resample-channel}
Let \(P^\star_\delta\) be the optimal coupling, represented by parameters \((u,\kappa)\).
Denote
\[
A=\sum_x u_x,\qquad \pi(x):=\frac{u_x}{A},
\qquad
p_x:=\mathbb P_{P^\star_\delta}(Y=x\mid X=x)=\frac{(1+\kappa)u_x}{A+\kappa u_x}\in(0,1).
\]
Then the conditional distribution \(Y\mid X=x\) has the decomposition
\[
\boxed{\;
\mathbb P(Y=x\mid X=x)=p_x,\qquad
\mathbb P(Y=y\mid X=x,\ y\neq x)=(1-p_x)\frac{\pi(y)}{1-\pi(x)}\; }.
\]
\end{theorem}
\begin{proof}
By the exponential family form, for fixed \(x\) we have
\(
P^\star_\delta(x,x)\propto (1+\kappa)u_x^2
\)
while for \(y\neq x\) we have
\(
P^\star_\delta(x,y)\propto u_xu_y
\).
Normalizing over \(y\) and writing in terms of \(\pi\) yields the decomposition. This completes the proof.
\end{proof}

\begin{theorem}[Mutual Information Characterization of Covering Threshold: Minimum Mutual Information under Diagonal Constraint]\label{thm:pom-typeclass-rate-distortion-mutual-information}
Let
\[
\mathcal C_\delta(w):=\Bigl\{P\in\Delta(X\times X):\ P_X=P_Y=w,\ \mathbb P_P(X=Y)\ge 1-\delta\Bigr\}.
\]
Then the strict equality holds
\[
\boxed{\;
R_w(\delta)=\inf_{P\in\mathcal C_\delta(w)} I_P(X;Y)\; } ,
\]
and this infimum is attained by the unique coupling \(P^\star_\delta\).
\end{theorem}
\begin{proof}
For any \(P\in\mathcal C_\delta(w)\), the identity holds
\(
I_P(X;Y)=H(w)-H_P(Y|X)
\).
Hence minimizing mutual information over \(\mathcal C_\delta(w)\) is equivalent to maximizing conditional entropy.
By Theorem \ref{thm:pom-typeclass-hamming-ball-volume-variational},
the maximum conditional entropy value is \(\mathcal V_w(\delta)\), attained by the unique maximal coupling of Theorem \ref{thm:pom-typeclass-hamming-ball-unique-maximizer-exponential-family}.
Thus
\(
R_w(\delta)=H(w)-\mathcal V_w(\delta)=I_{P^\star_\delta}(X;Y)
\).
This completes the proof.
\end{proof}
\begin{theorem}[Dual Derivative Identity: \texorpdfstring{$R_w'(\delta)$}{Rw'(delta)} Equals Negative of Optimal Diagonal Multiplier]\label{thm:pom-typeclass-dual-derivative}
Write \(p_0=\sum_{x\in X}w(x)^2\).
When \(1-\delta>p_0\) (diagonal mass constraint active), the function \(\delta\mapsto R_w(\delta)\) is differentiable on this interval, satisfying
\[
\boxed{\;
-\frac{d}{d\delta}R_w(\delta)=\lambda(\delta)\; } ,
\]
where \(\lambda(\delta)\) is the KKT multiplier (shadow price of diagonal mass constraint) of the optimal coupling in Theorem \ref{thm:pom-typeclass-hamming-ball-unique-maximizer-exponential-family}.
When \(1-\delta\le p_0\), the optimal coupling is independent, \(R_w(\delta)=0\) and \(\lambda(\delta)=0\).
\end{theorem}
\begin{proof}
Write the diagonal mass constraint as
\(
g_\delta(P):=(1-\delta)-\mathbb P_P(X=Y)\le 0
\).
For the maximization problem
\(
\sup\{H_P(Y|X):\ P_X=P_Y=w,\ g_\delta(P)\le 0\}
\),
take Lagrangian
\[
\mathscr L(P,\lambda)
:=H_P(Y|X)-\lambda\,g_\delta(P)
=H_P(Y|X)+\lambda\bigl(\mathbb P_P(X=Y)-1+\delta\bigr),
\qquad \lambda\ge 0.
\]
By Theorem \ref{thm:pom-typeclass-hamming-ball-unique-maximizer-exponential-family}:
when \(1-\delta>p_0\) the constraint is active and saturates (\(\lambda(\delta)>0\));
when \(1-\delta\le p_0\) the optimal solution is independent coupling with \(\lambda(\delta)=0\).
In both cases strict concavity yields strong duality and uniqueness of multiplier.
Hence the optimal value can be written as
\[
\mathcal V_w(\delta)=\max_{P_X=P_Y=w}\ \mathscr L(P,\lambda(\delta)),
\]
with \(\lambda(\delta)\) consistent with the diagonal coefficient in the exponential family form.
By the envelope theorem,
\[
\frac{d}{d\delta}\mathcal V_w(\delta)=\frac{\partial}{\partial\delta}\mathscr L(P^\star_\delta,\lambda(\delta))=\lambda(\delta).
\]
Since \(R_w(\delta)=H(w)-\mathcal V_w(\delta)\), we have \(R_w'(\delta)=-\lambda(\delta)\), yielding the stated identity. This completes the proof.
\end{proof}
\endinput





% (User input window)
% 2026-02-14 10:34:19 (Asia/Singapore)

\paragraph{Variational Structure of Two-Temperature Failure Exponent: Power-Law Coupling, Majorization Chain, and Critical Quadratic Law}
This paragraph lifts the KL-symmetric rate function of Theorem \ref{thm:pom-microcanonical-hypergeometric-ldp-kl-sym}
from the single-variable formulation \((q\mapsto J_\beta(q))\) to a two-temperature variational problem, and gives the structure of the unique minimizer on the supercritical side of the entropy threshold.
The result serves as an exponential-level failure certificate for the "residual composition entropy too large" event, compatible with the critical window narrative of Theorem \ref{thm:pom-microcanonical-query-plus-bits-phase}.

\begin{definition}[Two-Temperature Failure Exponent: Minimum KL-Symmetric Cost for Entropy Threshold]\label{def:pom-microcanonical-two-temperature-failure-exponent}
Let \(X\) be a finite set (\(|X|=n\ge 2\)), take full-support distribution \(w\in\Delta(X)\), and fix \(\beta\in(0,1)\).
For any \(q,r\in\Delta(X)\) satisfying mass conservation
\[
w=\beta q+(1-\beta)r,
\]
define the two-temperature divergence functional
\[
J_\beta(q,r):=\beta D(q\|w)+(1-\beta)D(r\|w).
\]
For entropy threshold \(h>H(w)\), define the supercritical failure exponent
\[
E_{\mathrm{fail}}(\beta,h)
:=
\inf\Bigl\{J_\beta(q,r):\ q,r\in\Delta(X),\ w=\beta q+(1-\beta)r,\ H(r)\ge h\Bigr\}.
\]
If the feasible domain is empty, we conventionalize \(E_{\mathrm{fail}}(\beta,h)=+\infty\).
\end{definition}

\begin{theorem}[Two-Temperature KKT Structure: Power-Law Coupling of Unique Minimizer]\label{thm:pom-microcanonical-two-temperature-kkt-power-law}
Following Definition \ref{def:pom-microcanonical-two-temperature-failure-exponent}, assume given \(h\) such that the feasible domain is non-empty and \(h>\!H(w)\).
Then the minimization problem has a unique minimum point \((q^\star,r^\star)\), with entropy constraint saturating:
\[
H(r^\star)=h.
\]
Moreover, there exists unique multiplier \(\eta^\star>0\) (hence unique exponent \(\gamma>1\)) such that for all \(x\in X\),
\[
q^\star(x)=\frac{(r^\star(x))^{\gamma}}{\sum_{y\in X}(r^\star(y))^{\gamma}},
\qquad
\gamma=1+\frac{\eta^\star}{1-\beta},
\]
equivalently
\(
\log q^\star=\gamma\log r^\star+\mathrm{const}
\)
(constant independent of \(x\)).
\end{theorem}
\begin{proof}
The feasible domain is given by linear equality and upper level set of concave function \(H\), hence is compact and convex; while \(J_\beta\) as a positively weighted sum of two KL divergences is strictly convex in the interior of \(\Delta(X)^2\), yielding unique minimum.
\par
Since \(w\) has full support and \(\beta\in(0,1)\), if \((q,r)\) is feasible, then \(q,r\) automatically have full support; hence we can write the KKT system in the interior.
If the entropy constraint is not saturated, the corresponding multiplier \(\eta^\star=0\), hence the minimum must also minimize the relaxed problem containing only mass conservation; but the unique minimum of that relaxed problem is \(q=r=w\), contradicting \(h>H(w)\).
Thus \(\eta^\star>0\) and \(H(r^\star)=h\).
\par
Writing the Lagrangian for the constraint and taking stationarity condition for each coordinate, eliminating the mass conservation multiplier yields
\(
\log r^\star=\alpha\log q^\star+\mathrm{const}
\),
where
\(
\alpha=\frac{1-\beta}{1-\beta+\eta^\star}\in(0,1)
\).
Hence \(q^\star\propto (r^\star)^{1/\alpha}\), setting \(\gamma:=1/\alpha=1+\eta^\star/(1-\beta)\) yields the stated power-law coupling.
Uniqueness forces \(\eta^\star\) and \(\gamma\) to be unique. This completes the proof.
\end{proof}

\begin{proposition}[Existence, Uniqueness, and Strict Monotonicity of Entropy along Power-Law Orbit]\label{prop:pom-microcanonical-two-temperature-escort-orbit-entropy-monotone}
Following Definition \ref{def:pom-microcanonical-two-temperature-failure-exponent}.
For each \(\gamma>1\), define power escort transform
\[
\mathrm{Escort}_\gamma(r)(x):=\frac{r(x)^\gamma}{\sum_{y\in X}r(y)^\gamma}\qquad(r\in\Delta(X)).
\]
If there exist \((q,r)\in\Delta(X)^2\) satisfying \(w=\beta q+(1-\beta)r\) and \(q=\mathrm{Escort}_\gamma(r)\), then this solution is unique; denote it \((q_\gamma,r_\gamma)\).
Moreover, on its domain of definition, the map
\[
\gamma\longmapsto H(r_\gamma)
\]
is continuous and strictly increasing, satisfying
\[
\lim_{\gamma\downarrow 1}r_\gamma=w.
\]
Let
\[
H_{\max}(\beta,w):=\sup\Bigl\{H(r):\ r\in\Delta(X),\ \frac{w-(1-\beta)r}{\beta}\in\Delta(X)\Bigr\}\in(H(w),\log n],
\]
then \(\lim_{\gamma\to\infty}H(r_\gamma)=H_{\max}(\beta,w)\).
In particular, if the uniform distribution \(U\) is feasible (equivalent to \(w(x)\ge (1-\beta)/n\) for all \(x\)), then \(H_{\max}(\beta,w)=\log n\) and \(r_\gamma\to U\).
Therefore for any \(h\in(H(w),H_{\max}(\beta,w))\) there exists unique \(\gamma(h)>1\) such that \(H(r_{\gamma(h)})=h\), and this \((q_{\gamma(h)},r_{\gamma(h)})\) coincides with the minimizer of Theorem \ref{thm:pom-microcanonical-two-temperature-kkt-power-law}.
\end{proposition}
\begin{proof}
For fixed \(\eta>0\) consider the penalized problem
\[
\mathcal{V}(\eta):=\inf\Bigl\{J_\beta(q,r)-\eta H(r):\ q,r\in\Delta(X),\ w=\beta q+(1-\beta)r\Bigr\}.
\]
Since \(-H\) is strictly convex and \(J_\beta\) is strictly convex, this problem has a unique minimum \((q_\eta,r_\eta)\);
its KKT system yields the power-law relation \(q_\eta=\mathrm{Escort}_\gamma(r_\eta)\) with \(\gamma=1+\eta/(1-\beta)>1\).
Hence for each \(\gamma>1\) there is at most one feasible solution, denoted \((q_\gamma,r_\gamma)\).
\par
The function \(\eta\mapsto \mathcal{V}(\eta)\) is the infimum envelope in \(\eta\), hence concave; and at the unique minimum is differentiable satisfying
\(
\mathcal{V}'(\eta)=-H(r_\eta)
\).
Concavity implies \(-\mathcal{V}'(\eta)=H(r_\eta)\) is strictly increasing in \(\eta\) (hence in \(\gamma\)), continuity follows from implicit function/stability and unique minimality.
\par
As \(\eta\downarrow 0\), the penalty term vanishes and the unique minimum of the relaxed problem is \(q=r=w\), hence \(r_\gamma\to w\) (\(\gamma\downarrow 1\)).
As \(\eta\to\infty\), \(-\eta H(r)\) dominates and forces \(r_\eta\) to converge to the unique solution of "maximizing \(H\) within the feasible polyhedron", hence \(H(r_\gamma)\uparrow H_{\max}(\beta,w)\);
if \(U\) is feasible, the maximum entropy point is \(U\) with limit \(\log n\).
Finally, for \(h\in(H(w),H_{\max})\) strict monotonicity yields unique \(\gamma(h)\), matching the unique minimizer of Theorem \ref{thm:pom-microcanonical-two-temperature-kkt-power-law}. This completes the proof.
\end{proof}

\begin{corollary}[Common Ordering and Majorization Chain: \texorpdfstring{$r^\star\prec w\prec q^\star$}{r* ≺ w ≺ q*}]\label{cor:pom-microcanonical-two-temperature-majorization-chain}
In the setting of Theorem \ref{thm:pom-microcanonical-two-temperature-kkt-power-law}, there exists a single permutation such that the coordinates of \(r^\star,w,q^\star\) arranged in non-increasing order satisfy the cumulative sum chain
\[
\sum_{i=1}^k r^\star_{(i)}\ \le\ \sum_{i=1}^k w_{(i)}\ \le\ \sum_{i=1}^k q^\star_{(i)}\qquad(k=1,\dots,n-1),
\]
with at least one \(k\) strict inequality; hence \(r^\star\prec w\prec q^\star\).
\end{corollary}
\begin{proof}
By power-law coupling \(q^\star=\mathrm{Escort}_\gamma(r^\star)\) (\(\gamma>1\)), they have the same coordinate ordering, and power transformation strictly cools the distribution in majorization sense: \(r^\star\prec q^\star\) (Hardy-Littlewood-Pólya).
Since \(w=\beta q^\star+(1-\beta)r^\star\) is a convex combination of vectors with the same ordering, cumulative sums are linear in convex combinations, yielding the chain inequality and \(r^\star\prec w\prec q^\star\).
Strictness follows from \(h>H(w)\) excluding \(r^\star=w\) (hence also excluding \(q^\star=w\)). This completes the proof.
\end{proof}

\begin{proposition}[Jensen-Shannon Structure Identity: \texorpdfstring{$J_\beta$}{Jbeta} Equals Mixture Entropy Deficit]\label{prop:pom-microcanonical-two-temperature-js-identity}
Under constraint \(w=\beta q+(1-\beta)r\), for any feasible \((q,r)\) the identity holds
\[
J_\beta(q,r)=H(w)-\beta H(q)-(1-\beta)H(r).
\]
Hence when the feasible domain is non-empty and \(h>H(w)\),
\[
E_{\mathrm{fail}}(\beta,h)
=H(w)-(1-\beta)h-\beta\sup\Bigl\{H(q):\exists r,\ w=\beta q+(1-\beta)r,\ H(r)=h\Bigr\}.
\]
\end{proposition}
\begin{proof}
By \(D(p\|w)=-H(p)-\sum_x p(x)\log w(x)\), and using mass conservation
\(
\beta q+(1-\beta)r=w
\)
to cancel \(\sum_x(\beta q+(1-\beta)r)\log w\), we obtain the first formula.
The second formula substitutes \(H(r^\star)=h\) from Theorem \ref{thm:pom-microcanonical-two-temperature-kkt-power-law} and performs equivalent maximization over \(q\). This completes the proof.
\end{proof}
\begin{theorem}[Strong Duality, Unique Multiplier, and Sensitivity Formula]\label{thm:pom-microcanonical-two-temperature-strong-duality-sensitivity}
Following Definition \ref{def:pom-microcanonical-two-temperature-failure-exponent}, assume given \(h\in(H(w),H_{\max}(\beta,w))\) such that the feasible domain is non-empty.
Define dual value function (\(\eta\ge 0\))
\[
\mathcal{V}(\eta):=\inf\Bigl\{J_\beta(q,r)-\eta H(r):\ q,r\in\Delta(X),\ w=\beta q+(1-\beta)r\Bigr\}.
\]
Then strong duality holds
\[
E_{\mathrm{fail}}(\beta,h)=\sup_{\eta\ge 0}\bigl(\mathcal{V}(\eta)+\eta h\bigr),
\]
with supremum attained at unique \(\eta^\star>0\).
Moreover, the function \(h\mapsto E_{\mathrm{fail}}(\beta,h)\) is differentiable on \((H(w),H_{\max}(\beta,w))\), and
\[
\frac{\partial}{\partial h}E_{\mathrm{fail}}(\beta,h)=\eta^\star,
\qquad
\gamma(h)=1+\frac{\eta^\star}{1-\beta},
\]
where \(\gamma(h)\) is the unique power exponent of Theorem \ref{thm:pom-microcanonical-two-temperature-kkt-power-law}.
\end{theorem}
\begin{proof}
The constraint \(H(r)\ge h\) is equivalent to convex inequality \(h-H(r)\le 0\), and satisfies Slater condition when \(h< H_{\max}(\beta,w)\) (strict feasible point exists).
Hence strong duality holds, and since the primal minimizer is unique, the KKT multiplier is also unique satisfying \(\eta^\star>0\) (same as Theorem \ref{thm:pom-microcanonical-two-temperature-kkt-power-law}).
Differentiability and derivative formula follow from general sensitivity theorem (envelope theorem) of convex duality;
the last formula substitutes the power-law relation \(\gamma=1+\eta/(1-\beta)\). This completes the proof.
\end{proof}
\begin{theorem}[Critical Entropy Window Quadratic Law: Precise Jump-Off Coefficient as \texorpdfstring{$h\downarrow H(w)$}{h↓H(w)}]\label{thm:pom-microcanonical-two-temperature-critical-quadratic-law}
Assume \(w\) has full support and is non-uniform, write
\[
V(w):=\Var_{X\sim w}\bigl(\log w(X)\bigr)>0
\]
(same as Corollary \ref{cor:pom-microcanonical-information-clt}).
Let \(\delta:=h-H(w)\downarrow 0^+\) with \(h\in(H(w),H_{\max}(\beta,w))\).
Then asymptotic expansion holds
\[
E_{\mathrm{fail}}(\beta,h)
=\frac{1-\beta}{2\beta}\cdot \frac{\delta^2}{V(w)}+o(\delta^2).
\]
Equivalently, using bit density
\(
b:=\frac{(1-\beta)h}{\log 2}
\)
and critical point
\(
b_{\mathrm{crit}}:=\frac{(1-\beta)H(w)}{\log 2}
\),
\[
E_{\mathrm{fail}}(\beta,b)
=\frac{(b-b_{\mathrm{crit}})^2(\log 2)^2}{2\beta(1-\beta)V(w)}+o\!\left((b-b_{\mathrm{crit}})^2\right).
\]
\end{theorem}
\begin{proof}
As \(\delta\downarrow 0\), the optimal point \((q^\star,r^\star)\) approaches \((w,w)\).
Write \(r=w+u\) (\(\sum_x u(x)=0\)), then by mass conservation \(q=w-\frac{1-\beta}{\beta}u\).
Second-order expansion of KL at \(w\) yields
\[
J_\beta(q,r)
=\frac{1-\beta}{2\beta}\sum_{x\in X}\frac{u(x)^2}{w(x)}+o(\|u\|_2^2).
\]
On the other hand, first-order variation of entropy (using \(\sum_x u(x)=0\)) gives
\[
H(w+u)-H(w)=-\sum_{x\in X}u(x)\log w(x)+o(\|u\|_2),
\]
hence leading-order constraint is \(\sum_x u(x)\log w(x)=-\delta\).
Thus leading-order problem reduces to: minimize quadratic form \(\sum u^2/w\) under linear constraint \(\sum u\log w=-\delta\).
Lagrange multiplier method yields unique leading direction
\[
u^\star(x)=-\frac{\delta}{V(w)}\,w(x)\Bigl(\log w(x)-\EE_w[\log w]\Bigr),
\]
satisfying
\(
\sum_x\frac{(u^\star(x))^2}{w(x)}=\delta^2/V(w)
\).
Substituting back yields
\(
E_{\mathrm{fail}}(\beta,h)=\frac{1-\beta}{2\beta}\cdot \frac{\delta^2}{V(w)}+o(\delta^2)
\).
Bit density form follows from substitution \(\delta=\frac{\log 2}{1-\beta}(b-b_{\mathrm{crit}})\). This completes the proof.
\end{proof}





\input{sections/body/pom/parts/part__pom-microcanonical-query-distortion-strong-converse}
\endinput








